\section{Specification and Design}
\label{sec:design}

\begin{guidance}[Specification and design (15 marks)]
Required (brief, A.3--A.5): an \textbf{embedded server FSM} covering IDLE,
IN\_PROGRESS, GAME\_OVER \textbf{and the move-validation branch}; a
\textbf{flowchart of the GUI application}; and a \textbf{circuit diagram} of the
STM-to-PC UART connection, with the on-board button and LED behaviour labelled.
Diagrams may be at function level. Explain \textbf{why} responsibilities sit on
the server or the client. Draw diagrams in a tool that exports vector PDF
(draw.io, TikZ, Mermaid) and check that every label is readable when printed.
\end{guidance}

\subsection{\ph{Heading about your overall architecture}}
\ph{How the firmware and the client are divided into modules, and why. Which
parts are pure logic, which touch hardware? Which side owns which decision
(board, turn, legality, win detection, redrawing)? A small table or block
diagram helps.}

\subsection{\ph{Heading about your server FSM}}
\ph{Walk the reader through Fig.~\ref{fig:fsm}: the three states, what each
command does in each state, and the move-validation branch (which checks, in
which order, and what happens when one fails). State which transitions the
figure omits, if any.}

\begin{figure*}[t]
\centering
\figph{5.5cm}{Server FSM (full width): states IDLE, IN\_PROGRESS, GAME\_OVER;
every transition labelled ``command [guard] / response''; the MOVE validation
branch drawn as a chain of checks, each with its NACK exit; the win/draw
outcome leading to GAME\_OVER; the button's effect.}
\caption{\ph{Server FSM with the move-validation branch.}}
\label{fig:fsm}
\end{figure*}

\subsection{\ph{Heading about your GUI}}
\ph{Walk through Fig.~\ref{fig:gui}: connecting, sending a move, handling each
kind of reply, redrawing, disabling moves after game over, disconnecting,
and what happens on a timeout or an unplugged cable. Explain how the GUI
avoids freezing while it waits for the board (threads, timers or polling).}

\begin{figure}[t]
\centering
\figph{7cm}{GUI flowchart: start $\rightarrow$ connect $\rightarrow$ wait for
input $\rightarrow$ send request $\rightarrow$ branch on the reply (board, win/draw,
NACK, timeout) $\rightarrow$ update display $\rightarrow$ loop; plus the
disconnect and error paths.}
\caption{\ph{GUI application flowchart.}}
\label{fig:gui}
\end{figure}

\subsection{\ph{Heading about your circuit}}
\ph{Describe Fig.~\ref{fig:circuit}: how USART TX/RX reach the PC (on-board
virtual COM port or a USB-UART converter, and which one), ground, the button
and LED pins you use, their active level, and what each LED and button does in
your game.}

\begin{figure}[t]
\centering
\figph{3.5cm}{Circuit diagram: STM32F051 TX/RX pins $\rightarrow$ USB-UART path
$\rightarrow$ PC, with GND; button and LED pins labelled with their behaviour.}
\caption{\ph{Circuit: UART link to the PC, button and LEDs.}}
\label{fig:circuit}
\end{figure}
